\subsection{A fixed-weight formulation of the simultaneous field problem}

\begin{proposition}[KT-JACOBI-FIXED-WEIGHT-REFORMULATION]
Let $p\ge31$ be prime with $p\equiv3\pmod4$, put
$b=(p+3)/2$, and let $D<0$ be an odd fundamental discriminant
prime to $p$.  Set
\[
 d=|D|,\qquad F=\mathbf Q(\sqrt{pd}),\qquad
 \chi_F=\chi_D\chi_{-p}.
\]
The last equality is an equality of primitive quadratic characters,
and $pd$ is the fundamental discriminant of $F$.  Both sides of
\[
 \frac{B_{b,\chi_D}}b\equiv\frac{B_{2,\chi_F}}2\pmod p
                                                        \tag{JFW.1}
\]
are $p$-integral.  Thus, for every
$e=(p+1)/2+(p-1)t$, $t\ge0$, the three auxiliary units are
equivalent to the fixed-weight conditions
\[
 p\nmid h(D),\qquad
 \frac{B_{2,\chi_F}}2\in\mathbf Z_p^\times,\qquad
 \frac{B_{3,\chi_D}}3\in\mathbf Z_p^\times.              \tag{JFW.2}
\]
If $d\le3p+6$, the first condition follows from
KT-SMALL-DISCRIMINANT-CLASS-BOUND.  This identifies the two
remaining special values at the same $D$; it does not assert their
simultaneous nonvanishing.
\end{proposition}

\begin{proof}
Since $D$ and $-p$ are coprime odd fundamental discriminants,
their product $pd$ is a positive fundamental discriminant.  Euler's
criterion gives, for every $a$ prime to $pd$,
\[
 \chi_D(a)a^{b-1}
 =\chi_D(a)a^{(p-1)/2+1}
 \equiv\chi_F(a)a\pmod p.                              \tag{JFW.3}
\]
Choose, by the Chinese remainder theorem, an integer $c>1$ with
$c\equiv1\pmod p$, $(c,d)=1$, and $\chi_D(c)=-1$.
Then $\chi_F(c)=-1$ as well.  Apply Guerzhoy's Proposition~2
with $\Delta=d$, $(\xi_1,r_1)=(\chi_D,b)$ and
$(\xi_2,r_2)=(\chi_F,2)$.  Its hypotheses allow the second
character to have conductor divisible by $p$.
The congruence of the integrands is exactly (JFW.3).  The two
regularizing factors are
\[
 1-\chi_D(c)c^b\equiv
 1-\chi_F(c)c^2\equiv2\pmod p,
\]
and hence are units.  The first Euler factor is
$1-\chi_D(p)p^{b-1}\equiv1\pmod p$; the second is exactly
$1$, since $\chi_F(p)=0$.

Proposition~1 expresses each regularized value as an integral
against an integral measure.  Thus both unregularized values
are integral, and Proposition~2 permits division by the common
unit residue $2$.  Using $L(1-r,\xi)=-B_{r,\xi}/r$ proves
(JFW.1).  In particular no prime-to-$p$ conductor assumption has
been inserted for $\chi_F$.

By KT-FIELD-CHARACTER-REDUCTION, the first twisted unit is
equivalent to $B_{b,\chi_D}/b$ being a unit, while
$2e\equiv2\pmod{p-1}$ identifies the second twisted unit
with $B_{3,\chi_D}/3$.  The untwisted $B_1$ unit is
equivalent to $p\nmid h(D)$; the factors from the roots of
unity are units because $p\ge31$.  This proves (JFW.2) in
every indicated weight.  Finally $3p+6<4p$ for $p\ge31$,
so the cited class-number bound applies.
\end{proof}

\paragraph{Exact primary input and scope.}
P.~Guerzhoy, \emph{Quadratic congruences for Cohen--Eisenstein
series}, Glasgow Math.~J.~41 (1999), 141--144,
\url{https://doi.org/10.1017/S001708959997060X}:
Proposition~1 and formula~(2), printed p.~143, give the integral
Mazur measure; Proposition~2, printed pp.~143--144,
and formula~(4) give the change-of-character Kummer congruence.
The short deduction of Proposition~2 from the measure is included
there.  The construction of that measure is an imported published
input, not a newly checked proof here.  Only these propositions
are used, not the final type-a/type-b congruence argument.
The local PDF and retrieval metadata are
\texttt{literature/jacobi\_simultaneous/Guerzhoy\_1999.pdf}
and its corresponding retrieval JSON.
This proposition is a reformulation of the simultaneous target,
not an additional vanishing result for the original rational class
component or an exclusion of genuine rank-three sources.
